\section{Giới thiệu: Người dùng, Bài toán và Lý do cần Agent}
\label{sec:intro}

\noindent \textbf{Người dùng mục tiêu và bài toán thực tế (\texttt{AT1}).} Trong vận hành ngân hàng bán lẻ, nhân viên nghiệp vụ thường cần tổng hợp nhanh số liệu giao dịch, tài khoản hoặc hồ sơ khoản vay mà không phụ thuộc vào đội ngũ kỹ thuật viết SQL. Trên lược đồ \texttt{financial.sqlite} (8 bảng, 61 cột), hệ thống xác định ba nhóm người dùng rõ ràng với phạm vi quyền tách biệt (không phân cấp cao--thấp tuyệt đối): (1) \textbf{\texttt{operator}} tra cứu hỗ trợ tài khoản, giao dịch và lệnh thanh toán định kỳ; (2) \textbf{\texttt{analyst}} chuyên viên BI chỉ được xuất báo cáo thống kê tổng hợp theo vùng/tháng, không được xem hồ sơ cá thể; và (3) \textbf{\texttt{credit\_officer}} theo dõi danh mục khoản vay sau giải ngân nhưng bị chặn thông tin thẻ và ghi chú dịch vụ. Ba kịch bản đại diện gồm: (i) truy vấn hợp lệ nhiều bảng kèm điều kiện thời gian/vùng; (ii) câu hỏi hỗn hợp yêu cầu cả cột hợp lệ lẫn cột định danh bị cấm (\texttt{id\_number}); và (iii) câu hỏi vô tình hoặc cố ý chạm dữ liệu chứa chỉ thị tiêm nhiễm trong cột \texttt{account.service\_note} hay \texttt{order.note}.

\noindent \textbf{Vì sao cần kiến trúc Agent thay vì script tĩnh hoặc single-prompt?} Một bộ quy tắc tĩnh không thể xử lý câu hỏi tự nhiên đa dạng từ người dùng, trong khi một lời gọi LLM đơn lẻ (\textit{single-pass}) dễ sai cú pháp hoặc nhầm tên bảng trên lược đồ lớn và thiếu cơ chế thu hẹp truy vấn khi gặp cột cấm. Ngược lại, một Agent thiếu kiểm soát lại tự làm suy yếu bảo mật nếu dựa vào system prompt để nhắc phân quyền. SQLGuard chọn mô hình tác tử có kiểm soát: LLM chỉ đóng vai trò đề xuất SQL, còn toàn bộ việc kiểm tra quyền, sửa lỗi có giới hạn (\textit{Bounded Reflection}) và làm sạch kết quả được thực thi xác định bằng mã nguồn.

\noindent \textbf{Đóng góp chính.} (1) Phòng thủ 3 lớp mức cột gồm \texttt{PolicySQLGate} (SQLGlot AST + \textit{Partial Answer} + ẩn nhóm $k<5$ cho \texttt{analyst}), \texttt{SecureSqlTool} và \texttt{SQLiteAuthorizer}; (2) Runtime MCP 2 công cụ tối giản với ngân sách thực thi chặt và kiểm toán băm; (3) Cô lập môi trường kiểm thử \textit{Attack Lab} cùng đối chứng trực tiếp \texttt{NaiveAgent} trên cùng Gemini LLM; và (4) Bộ đánh giá định lượng \textit{Security--Utility} cùng bộ tái lập một lệnh trên BIRD Financial.
